% ============================================================
%  IPPE for a square marker --- full derivation, every step
%  Standalone note. Compiles with pdfLaTeX (on Overleaf: just upload it).
%  Study aid kept alongside the code; it is not included in report.tex.
% ============================================================
\documentclass[11pt,a4paper]{article}

\usepackage[utf8]{inputenc}
\usepackage[T1]{fontenc}
\usepackage[english]{babel}
\usepackage[a4paper,margin=2.5cm]{geometry}
\usepackage{amsmath,amssymb,amsthm,mathtools}
\usepackage{booktabs,array,longtable}
\usepackage{enumitem}
\usepackage[colorlinks=true,linkcolor=blue!55!black]{hyperref}

\theoremstyle{remark}
\newtheorem*{remark}{Remark}

% --- notation shortcuts (none of these clash with LaTeX built-ins) ---
\newcommand{\Rset}{\mathbb{R}}
\newcommand{\Rmat}{\mathbf{R}}
\newcommand{\Rtil}{\widetilde{\mathbf{R}}}
\newcommand{\Amat}{\mathbf{A}}
\newcommand{\Bmat}{\mathbf{B}}
\newcommand{\Hmat}{\mathbf{H}}
\newcommand{\Jmat}{\mathbf{J}}
\newcommand{\Kmat}{\mathbf{K}}
\newcommand{\Lmat}{\mathbf{L}}
\newcommand{\Mmat}{\mathbf{M}}
\newcommand{\Rvee}{\mathbf{R}_{\vvec}}
\newcommand{\Pimat}{\boldsymbol{\Pi}_{\vvec}}
\newcommand{\Idm}{\mathbf{I}}
\newcommand{\bvec}{\mathbf{b}}
\newcommand{\cvec}{\mathbf{c}}
\newcommand{\evec}{\mathbf{e}}
\newcommand{\lvec}{\mathbf{l}}
\newcommand{\Pvec}{\mathbf{P}}
\newcommand{\qvec}{\mathbf{q}}
\newcommand{\rvec}{\mathbf{r}}
\newcommand{\tvec}{\mathbf{t}}
\newcommand{\uvec}{\mathbf{u}}
\newcommand{\vvec}{\mathbf{v}}
\newcommand{\wvec}{\mathbf{w}}
\newcommand{\xvec}{\mathbf{x}}
\newcommand{\Xvec}{\mathbf{X}}
\newcommand{\zerovec}{\mathbf{0}}
\newcommand{\dhat}{\hat{\mathbf{d}}}
\newcommand{\skewm}[1]{[#1]_{\times}}

\title{IPPE for a Square Marker\\[0.3em]
       \large Every formula, every symbol, in order}
\date{}

\begin{document}
\maketitle

\noindent
This note derives the two IPPE pose candidates from scratch. Nothing is assumed
beyond the pinhole camera model and basic linear algebra. Every symbol is
defined in Section~\ref{sec:symbols}, and again where it first appears.

\tableofcontents

% ============================================================
\section{Symbols}
\label{sec:symbols}
% ============================================================

\begin{longtable}{@{}lp{11.2cm}@{}}
\toprule
Symbol & Meaning \\
\midrule
\endfirsthead
\toprule
Symbol & Meaning \\
\midrule
\endhead
$\Xvec = (X,Y,0)^{\top}$ & a point of the marker, in the marker frame; the
  marker lies on the plane $Z=0$, so the third coordinate is always zero \\
$\Rmat \in \mathrm{SO}(3)$ & rotation of the pose: marker frame $\to$ camera
  frame \\
$\tvec = (t_1,t_2,t_3)^{\top}$ & translation of the pose; $t_3 > 0$ is the
  depth of the marker centre \\
$\rvec_1,\rvec_2,\rvec_3$ & the three columns of $\Rmat$ \\
$\Pvec = (P_1,P_2,P_3)^{\top}$ & the marker point in the camera frame,
  $\Pvec = \Rmat\Xvec + \tvec$ \\
$\xvec = (x_1,x_2)^{\top}$ & normalized image coordinates of $\Pvec$, i.e.\
  pixels with $\Kmat^{-1}$ already applied \\
$\Kmat$ & camera intrinsic matrix; plays no role after normalization \\
$\Hmat$ & homography from the marker plane to normalized image coordinates,
  $3\times3$, scaled so that $h_{33}=1$ \\
$h_{ij}$ & entry $(i,j)$ of $\Hmat$ \\
$\vvec = (v_1,v_2)^{\top}$ & normalized image point of the \emph{square
  centre}; equals $(h_{13},h_{23})^{\top}$ \\
$\Jmat$ & $2\times2$ derivative of the normalized projection with respect to
  $(X,Y)$, evaluated at the square centre \\
$\Pimat$ & $2\times3$ matrix appearing in the derivative of the perspective
  divide; see \eqref{eq:pidef} \\
$\Rvee$ & rotation taking $\evec_3$ onto $(v_1,v_2,1)^{\top}$; a change of
  frame \\
$\Lmat$ & the pose rotation seen in that frame, $\Rmat = \Rvee\Lmat$ \\
$\lvec_1,\lvec_2,\lvec_3$ & the three columns of $\Lmat$ \\
$\Rtil$ & upper-left $2\times2$ block of $\Lmat$; the part IPPE recovers first \\
$\bvec = (b_1,b_2)^{\top}$ & third row of the first two columns of $\Lmat$; the
  part the derivative does \emph{not} see, and the source of the ambiguity \\
$\cvec = (c_1,c_2)^{\top}$, $a$ & the third column of $\Lmat$, split as
  $\lvec_3 = (c_1,c_2,a)^{\top}$ \\
$\Amat$ & $2\times2$ matrix defined by $\Amat = \Bmat^{-1}\Jmat$; turns out to
  equal $\Rtil / t_3$ \\
$\Bmat$ & $2\times2$ matrix defined in \eqref{eq:Bdef} \\
$\gamma$ & largest singular value of $\Amat$; turns out to equal $1/t_3$ \\
$\sigma_1,\sigma_2$ & singular values of a matrix, ordered
  $\sigma_1 \ge \sigma_2$ \\
$\theta$ & viewing angle: the angle between the marker normal and the line from
  the marker centre to the camera centre; $\theta = 0$ is fronto-parallel \\
$\dhat$ & unit line of sight from the camera centre to the marker centre,
  $\dhat = \tvec/\|\tvec\|$ \\
$\qvec^{(i)}$ & camera ray of corner $i$, $\qvec^{(i)} =
  (x_1^{(i)},x_2^{(i)},1)^{\top}$ \\
$\skewm{\qvec}$ & skew-symmetric matrix with
  $\skewm{\qvec}\wvec = \qvec \times \wvec$ \\
$\Idm$, $\evec_3$ & identity matrix; third basis vector $(0,0,1)^{\top}$ \\
\bottomrule
\end{longtable}

% ============================================================
\section{Step 0 --- The projection, and why a homography appears}
% ============================================================

A marker point $\Xvec = (X,Y,0)^{\top}$ lands in the camera frame at
\begin{equation}
  \Pvec \;=\; \Rmat\Xvec + \tvec
  \;=\; X\,\rvec_1 + Y\,\rvec_2 + \tvec ,
  \label{eq:rigid}
\end{equation}
because the third coordinate of $\Xvec$ is zero and therefore multiplies
$\rvec_3$ away. The camera then divides by depth:
\begin{equation}
  \xvec \;=\; \pi(\Pvec) \;=\;
  \begin{pmatrix} P_1/P_3 \\[2pt] P_2/P_3 \end{pmatrix}.
  \label{eq:divide}
\end{equation}
Combining \eqref{eq:rigid} and \eqref{eq:divide} in homogeneous form,
\begin{equation}
  \begin{pmatrix} x_1 \\ x_2 \\ 1 \end{pmatrix}
  \;\sim\;
  \begin{bmatrix} \rvec_1 & \rvec_2 & \tvec \end{bmatrix}
  \begin{pmatrix} X \\ Y \\ 1 \end{pmatrix}
  \;=\; \Hmat \begin{pmatrix} X \\ Y \\ 1 \end{pmatrix},
  \label{eq:homography}
\end{equation}
where $\sim$ means ``equal up to a non-zero scale factor''. The whole geometry
of a planar marker is thus a single $3\times3$ matrix
$\Hmat = [\,\rvec_1\ \rvec_2\ \tvec\,]$, defined up to scale. Fixing the scale
by dividing by $h_{33} = t_3$ gives
\begin{equation}
  \Hmat \;=\; \frac{1}{t_3}
  \begin{bmatrix} \rvec_1 & \rvec_2 & \tvec \end{bmatrix},
  \qquad h_{33} = 1 .
  \label{eq:Hscaled}
\end{equation}

\paragraph{The idea behind IPPE.} One could try to read $\rvec_1$, $\rvec_2$
and $\tvec$ straight off $\Hmat$, but an estimated $\Hmat$ is noisy and its
first two columns are not exactly orthonormal, so that route needs a correction
step of its own. IPPE instead uses only \emph{local} information at one point of
the plane: the image of the square centre, and the first derivative of the map
there. Two quantities, and everything else follows.

% ============================================================
\section{Step 1 --- The two local quantities}
% ============================================================

\subsection{The image of the centre}

The square is centred at the origin of the marker frame, so evaluate
\eqref{eq:homography} at $(X,Y) = (0,0)$. Only the last column of $\Hmat$
survives:
\begin{equation}
  \vvec \;=\; \begin{pmatrix} h_{13} \\ h_{23} \end{pmatrix}
  \;\overset{\eqref{eq:Hscaled}}{=}\;
  \begin{pmatrix} t_1/t_3 \\[2pt] t_2/t_3 \end{pmatrix}.
  \label{eq:vdef}
\end{equation}
So $\vvec$ is the normalized projection of the marker centre, obtained for free
from the third column of the homography.

\subsection{The derivative at the centre}

Differentiating the perspective divide \eqref{eq:divide} with respect to
$\Pvec$,
\begin{equation}
  \frac{\partial \pi}{\partial \Pvec}
  \;=\; \frac{1}{P_3}
  \begin{bmatrix}
    1 & 0 & -P_1/P_3 \\
    0 & 1 & -P_2/P_3
  \end{bmatrix}
  \;=\; \frac{1}{P_3}
  \begin{bmatrix}
    1 & 0 & -x_1 \\
    0 & 1 & -x_2
  \end{bmatrix},
  \label{eq:dpi}
\end{equation}
while \eqref{eq:rigid} shows that the derivative of $\Pvec$ with respect to the
plane coordinates is constant,
\begin{equation}
  \frac{\partial \Pvec}{\partial (X,Y)}
  \;=\; \begin{bmatrix} \rvec_1 & \rvec_2 \end{bmatrix}
  \qquad (3\times2).
  \label{eq:dP}
\end{equation}
At the centre $\Pvec = \tvec$, hence $P_3 = t_3$ and $\xvec = \vvec$. Writing
\begin{equation}
  \Pimat \;\coloneqq\;
  \begin{bmatrix}
    1 & 0 & -v_1 \\
    0 & 1 & -v_2
  \end{bmatrix}
  \qquad (2\times3),
  \label{eq:pidef}
\end{equation}
the chain rule applied to \eqref{eq:dpi} and \eqref{eq:dP} gives the second
local quantity,
\begin{equation}
  \boxed{\;
  \Jmat \;=\; \frac{1}{t_3}\,\Pimat
  \begin{bmatrix} \rvec_1 & \rvec_2 \end{bmatrix}
  \;}
  \qquad (2\times2).
  \label{eq:Jdef}
\end{equation}

\paragraph{How $\Jmat$ is computed in practice.} We do not know
$\rvec_1,\rvec_2,\tvec$; we know $\Hmat$. Writing the first normalized
coordinate as a ratio,
\begin{equation}
  x_1(X,Y) \;=\;
  \frac{h_{11}X + h_{12}Y + h_{13}}{h_{31}X + h_{32}Y + 1},
\end{equation}
the quotient rule at $(X,Y) = (0,0)$, where the denominator equals $1$, gives
$\partial x_1 / \partial X = h_{11} - h_{13}h_{31}$, and similarly for the other
three entries:
\begin{equation}
  \Jmat \;=\;
  \begin{bmatrix}
    h_{11} - h_{13}h_{31} & h_{12} - h_{13}h_{32} \\
    h_{21} - h_{23}h_{31} & h_{22} - h_{23}h_{32}
  \end{bmatrix}.
  \label{eq:JfromH}
\end{equation}
Equations \eqref{eq:Jdef} and \eqref{eq:JfromH} are the same matrix: the first
says what it \emph{means}, the second how to \emph{compute} it.

% ============================================================
\section{Step 2 --- A change of frame that removes one unknown}
% ============================================================

\subsection{The rotation $\Rvee$}

Let $\Rvee$ be any rotation with
\begin{equation}
  \Rvee\,\evec_3 \;=\;
  \frac{1}{\sqrt{v_1^2 + v_2^2 + 1}}
  \begin{pmatrix} v_1 \\ v_2 \\ 1 \end{pmatrix},
  \label{eq:Rvdef}
\end{equation}
that is, the rotation pointing the optical axis at the observed marker centre.
The minimal such rotation has a closed form through Rodrigues' formula; any
valid choice works. Define $\Lmat$ by
\begin{equation}
  \Rmat \;=\; \Rvee\,\Lmat
  \qquad \Longleftrightarrow \qquad
  \Lmat \;=\; \Rvee^{\top}\Rmat .
  \label{eq:Ldef}
\end{equation}
$\Lmat$ is the pose rotation expressed in the frame where the marker centre lies
on the optical axis. Since $\Rvee$ is known from $\vvec$ alone, recovering
$\Lmat$ is equivalent to recovering $\Rmat$.

\subsection{Substituting into $\Jmat$}

The first two columns of $\Rmat$ are
$[\,\rvec_1\ \rvec_2\,] = \Rvee[\,\lvec_1\ \lvec_2\,]$, so \eqref{eq:Jdef}
becomes
\begin{equation}
  \Jmat \;=\; \frac{1}{t_3}\,\Pimat\,\Rvee
  \begin{bmatrix} \lvec_1 & \lvec_2 \end{bmatrix}.
  \label{eq:Jsub}
\end{equation}
Split the $3\times2$ matrix $[\,\lvec_1\ \lvec_2\,]$ into its first two rows and
its third row:
\begin{equation}
  \begin{bmatrix} \lvec_1 & \lvec_2 \end{bmatrix}
  \;=\;
  \begin{bmatrix} \Rtil \\[2pt] \bvec^{\top} \end{bmatrix},
  \qquad
  \Rtil \coloneqq
  \begin{bmatrix}
    \ell_{11} & \ell_{12} \\
    \ell_{21} & \ell_{22}
  \end{bmatrix},
  \qquad
  \bvec \coloneqq
  \begin{pmatrix} \ell_{31} \\ \ell_{32} \end{pmatrix},
  \label{eq:split}
\end{equation}
where $\ell_{ij}$ is the entry $(i,j)$ of $\Lmat$. The product then separates
into two terms,
\begin{equation}
  \Pimat\Rvee
  \begin{bmatrix} \lvec_1 & \lvec_2 \end{bmatrix}
  \;=\;
  \underbrace{\big(\Pimat\,\Rvee^{(1,2)}\big)}_{\textstyle \eqqcolon\, \Bmat}
  \,\Rtil
  \;+\;
  \big(\Pimat\,\Rvee\,\evec_3\big)\,\bvec^{\top},
  \label{eq:twoterms}
\end{equation}
in which $\Rvee^{(1,2)}$ denotes the first two columns of $\Rvee$, so that
\begin{equation}
  \Bmat \;\coloneqq\; \Pimat\,\Rvee^{(1,2)}
  \qquad (2\times2).
  \label{eq:Bdef}
\end{equation}

\subsection{The cancellation}

The second term of \eqref{eq:twoterms} vanishes identically. By
\eqref{eq:Rvdef} the vector $\Rvee\evec_3$ is proportional to
$(v_1,v_2,1)^{\top}$, and by the definition \eqref{eq:pidef} of $\Pimat$,
\begin{equation}
  \Pimat \begin{pmatrix} v_1 \\ v_2 \\ 1 \end{pmatrix}
  \;=\;
  \begin{pmatrix} v_1 - v_1 \\ v_2 - v_2 \end{pmatrix}
  \;=\; \zerovec .
  \label{eq:cancel}
\end{equation}
This is the crucial step, and it deserves to be said in words: \emph{the
derivative of the image at the centre carries no information whatsoever about
$\bvec$.} Substituting \eqref{eq:cancel} into \eqref{eq:twoterms} and the result
into \eqref{eq:Jsub},
\begin{equation}
  \Jmat \;=\; \frac{1}{t_3}\,\Bmat\,\Rtil .
  \label{eq:JBR}
\end{equation}
Since $\Bmat$ is built from known quantities ($\vvec$, hence $\Pimat$ and
$\Rvee$), it can be inverted, which defines
\begin{equation}
  \boxed{\;
  \Amat \;\coloneqq\; \Bmat^{-1}\Jmat \;=\; \frac{1}{t_3}\,\Rtil
  \;}
  \label{eq:Adef}
\end{equation}
One $2\times2$ inversion therefore delivers the upper-left block of the
rotation, up to the unknown depth $t_3$.

% ============================================================
\section{Step 3 --- Fixing the scale}
% ============================================================

To turn $\Amat$ into $\Rtil$ we still need $t_3$. Orthonormality supplies it.

\subsection{The key identity}

The first two columns of $\Lmat$ are orthonormal, because $\Lmat$ is a rotation:
\begin{equation}
  \Mmat \;\coloneqq\;
  \begin{bmatrix} \lvec_1 & \lvec_2 \end{bmatrix}
  \;=\;
  \begin{bmatrix} \Rtil \\[2pt] \bvec^{\top} \end{bmatrix}
  \qquad \text{satisfies} \qquad
  \Mmat^{\top}\Mmat \;=\; \Idm_2 .
\end{equation}
Expanding the product block by block,
\begin{equation}
  \Mmat^{\top}\Mmat
  \;=\; \Rtil^{\top}\Rtil + \bvec\,\bvec^{\top}
  \;=\; \Idm_2 .
  \label{eq:key}
\end{equation}
Call this the \emph{key identity}. It is used twice: here, to fix the scale, and
in Section~\ref{sec:b}, to produce $\bvec$.

\subsection{Consequence for the singular values}

From \eqref{eq:key}, $\Rtil^{\top}\Rtil = \Idm_2 - \bvec\bvec^{\top}$. The
matrix $\bvec\bvec^{\top}$ has eigenvalue $\|\bvec\|^2$ along $\bvec$ and $0$ in
the orthogonal direction, so $\Rtil^{\top}\Rtil$ has eigenvalues $1$ and
$1 - \|\bvec\|^2$. Taking square roots, the singular values of $\Rtil$ are
\begin{equation}
  \sigma_1(\Rtil) = 1,
  \qquad
  \sigma_2(\Rtil) = \sqrt{1 - \|\bvec\|^2}.
  \label{eq:sv}
\end{equation}
The first is \emph{exactly} $1$, for every pose. Applying this to
\eqref{eq:Adef}, and using the fact that scaling a matrix scales its singular
values,
\begin{equation}
  \gamma \;\coloneqq\; \sigma_1(\Amat)
  \;=\; \frac{\sigma_1(\Rtil)}{t_3}
  \;=\; \frac{1}{t_3},
  \label{eq:gamma}
\end{equation}
and therefore
\begin{equation}
  \boxed{\;
  \Rtil \;=\; \frac{\Amat}{\gamma}
  \;}
  \qquad \text{and} \qquad
  t_3 \;=\; \frac{1}{\gamma}.
  \label{eq:Rtilde}
\end{equation}
The depth comes out as a by-product.

% ============================================================
\section{Step 4 --- Where $\bvec$ comes from}
\label{sec:b}
% ============================================================

\subsection{The equation for $\bvec$}

Rearranging the key identity \eqref{eq:key} isolates the unknown:
\begin{equation}
  \boxed{\;
  \bvec\,\bvec^{\top} \;=\; \Idm_2 - \Rtil^{\top}\Rtil
  \;}
  \label{eq:bb}
\end{equation}
Everything on the right is known once $\Rtil$ has been computed in
\eqref{eq:Rtilde}. That right-hand side is symmetric, positive semidefinite and
of rank at most one, exactly as an outer product $\bvec\bvec^{\top}$ must be. So
$\bvec$ is read off its leading eigenpair: if $\lambda \ge 0$ is the largest
eigenvalue and $\uvec$ a corresponding unit eigenvector, then
\begin{equation}
  \bvec \;=\; \pm\sqrt{\lambda}\;\uvec .
  \label{eq:beig}
\end{equation}

\subsection{The same thing, entry by entry}

Equation \eqref{eq:bb} says no more than that the columns of $\Mmat$ have unit
norm and are mutually orthogonal. Writing the entries of $\Rtil$ as
$\tilde r_{ij}$, the three scalar equations are
\begin{align}
  \tilde r_{11}^2 + \tilde r_{21}^2 + b_1^2 &= 1
  &&\Longrightarrow&
  b_1 &= \pm\sqrt{1 - \tilde r_{11}^2 - \tilde r_{21}^2},
  \label{eq:b1} \\
  \tilde r_{12}^2 + \tilde r_{22}^2 + b_2^2 &= 1
  &&\Longrightarrow&
  b_2 &= \pm\sqrt{1 - \tilde r_{12}^2 - \tilde r_{22}^2},
  \label{eq:b2} \\
  \tilde r_{11}\tilde r_{12} + \tilde r_{21}\tilde r_{22} + b_1 b_2 &= 0
  &&\Longrightarrow&
  \operatorname{sign}(b_1 b_2)
  &= -\operatorname{sign}\big(
     \tilde r_{11}\tilde r_{12} + \tilde r_{21}\tilde r_{22}\big).
  \label{eq:b3}
\end{align}
The first two equations give the magnitudes, the third ties the two signs
together. One global sign remains free, and that is the ambiguity.

\subsection{Why this is precisely the ambiguity}

Equation \eqref{eq:bb} is quadratic in $\bvec$: if $\bvec$ solves it, so does
$-\bvec$, since $(-\bvec)(-\bvec)^{\top} = \bvec\bvec^{\top}$. No data can
distinguish the two, because by the cancellation \eqref{eq:cancel} the local
derivative $\Jmat$ never saw $\bvec$ at all. Hence \emph{two} poses rather than
one. This is not a numerical accident nor an arbitrary tie-break: it is the
structure of the problem.

% ============================================================
\section{Step 5 --- Completing the rotation}
% ============================================================

For each sign choice we now know the first two columns of $\Lmat$,
\begin{equation}
  \lvec_1 =
  \begin{pmatrix} \tilde r_{11} \\ \tilde r_{21} \\ b_1 \end{pmatrix},
  \qquad
  \lvec_2 =
  \begin{pmatrix} \tilde r_{12} \\ \tilde r_{22} \\ b_2 \end{pmatrix}.
\end{equation}
A rotation has orthonormal columns and determinant $+1$, so the third column is
forced:
\begin{equation}
  \lvec_3 \;=\; \lvec_1 \times \lvec_2
  \;\eqqcolon\;
  \begin{pmatrix} c_1 \\ c_2 \\ a \end{pmatrix}.
  \label{eq:third}
\end{equation}
Assembling $\Lmat$ and undoing the change of frame \eqref{eq:Ldef} gives the two
candidates
\begin{equation}
  \boxed{\;
  \Rmat^{\pm} \;=\; \Rvee
  \begin{bmatrix}
    \Rtil & \pm\cvec \\[2pt]
    \pm\bvec^{\top} & a
  \end{bmatrix}
  \;}
  \label{eq:twoposes}
\end{equation}
Note that $\cvec$ flips sign together with $\bvec$: replacing $\bvec$ by
$-\bvec$ in \eqref{eq:third} reverses the first two components of the cross
product, while the third,
$a = \tilde r_{11}\tilde r_{22} - \tilde r_{12}\tilde r_{21}$, is unchanged.

% ============================================================
\section{Step 6 --- Translation}
% ============================================================

Once a rotation is fixed, translation is a linear problem. Let corner $i$ have
normalized image point $(x_1^{(i)}, x_2^{(i)})$ and camera ray
\begin{equation}
  \qvec^{(i)} =
  \begin{pmatrix} x_1^{(i)} \\ x_2^{(i)} \\ 1 \end{pmatrix}.
\end{equation}
The camera-frame point $\Rmat\Xvec^{(i)} + \tvec$ must lie along that ray, that
is, be parallel to it, which is expressed by a vanishing cross product:
\begin{equation}
  \qvec^{(i)} \times \big(\Rmat\Xvec^{(i)} + \tvec\big) = \zerovec
  \qquad \Longleftrightarrow \qquad
  \skewm{\qvec^{(i)}}\,\tvec
  \;=\; -\skewm{\qvec^{(i)}}\,\Rmat\Xvec^{(i)},
  \label{eq:translation}
\end{equation}
where $\skewm{\qvec}$ is the skew-symmetric matrix satisfying
$\skewm{\qvec}\wvec = \qvec \times \wvec$,
\begin{equation}
  \skewm{\qvec} =
  \begin{bmatrix}
    0 & -q_3 & q_2 \\
    q_3 & 0 & -q_1 \\
    -q_2 & q_1 & 0
  \end{bmatrix}.
\end{equation}
Only two of the three rows per corner are independent. Writing
$\Pvec = \Rmat\Xvec^{(i)} + \tvec$, the three rows of \eqref{eq:translation} are
\begin{align}
  &\text{row 1:} & -P_2 + x_2 P_3 &= 0, \\
  &\text{row 2:} & P_1 - x_1 P_3 &= 0, \\
  &\text{row 3:} & -x_2 P_1 + x_1 P_2 &= 0,
\end{align}
and row 3 equals $-(x_1\cdot\text{row 1} + x_2\cdot\text{row 2})$. It carries no
new information, but in a least-squares fit it would still change the weight of
each residual. The implementation keeps rows 1 and 2 only: stacked over the four
corners they give an $8 \times 3$ system, solved in the least-squares sense,
which already has rank $3$ with two corners.

\begin{remark}
Rows 1 and 2 are exactly the equations OpenCV uses, obtained by clearing the
denominator in $x_1 = (\rvec_1^{\top}\Xvec + t_1)/(\rvec_3^{\top}\Xvec + t_3)$
and its counterpart for $x_2$. Keeping all three rows instead gives a
$12\times3$ system with the same solution whenever the rotation fits the data
exactly, but a slightly different one otherwise, as for the mirror candidate.
With two rows the two implementations agree on both candidates.
\end{remark}

% ============================================================
\section{What the quantities mean geometrically}
\label{sec:geometry}
% ============================================================

The quantities $\gamma$, $\Rtil$, $\bvec$ and $a$ were introduced as algebraic
objects. Each of them turns out to have a direct geometric meaning, and the one
that matters most is $\bvec$.

\subsection{The line of sight and the viewing angle}

Let
\begin{equation}
  \dhat \;=\; \frac{\tvec}{\|\tvec\|}
  \label{eq:dhat}
\end{equation}
be the unit vector along the line of sight, from the camera centre to the marker
centre, expressed in the camera frame. The marker normal in the camera frame is
the third column $\rvec_3$ of $\Rmat$. The viewing angle $\theta$ is the angle
between that normal and the direction from the marker back to the camera,
$-\dhat$; for a marker facing the camera, $0 \le \theta < 90^{\circ}$. Hence
\begin{equation}
  \dhat \cdot \rvec_3 \;=\; -\cos\theta .
  \label{eq:dn}
\end{equation}

\subsection{The theorem: $\|\bvec\| = \sin\theta$}

\paragraph{Step 1 --- $\bvec$ as dot products.} By \eqref{eq:Ldef},
$\Lmat = \Rvee^{\top}\Rmat$, and by \eqref{eq:split}, $\bvec$ is the third row of
the first two columns of $\Lmat$. The entry $(3,j)$ of a product
$\Rvee^{\top}\Rmat$ is the third column of $\Rvee$ dotted with the $j$-th column
of $\Rmat$:
\begin{equation}
  b_j \;=\; \big(\Rvee\,\evec_3\big) \cdot \rvec_j ,
  \qquad j = 1,2 .
  \label{eq:bdot}
\end{equation}

\paragraph{Step 2 --- $\Rvee\evec_3$ is the line of sight.} By \eqref{eq:Rvdef},
$\Rvee\evec_3$ is the unit vector along $(v_1,v_2,1)^{\top}$, and by
\eqref{eq:vdef}, $(v_1,v_2,1)^{\top} = \tvec / t_3$. A unit vector along
$\tvec$ is $\dhat$, so
\begin{equation}
  \Rvee\,\evec_3 \;=\; \dhat
  \qquad\Longrightarrow\qquad
  \boxed{\;
  \bvec \;=\; \begin{pmatrix} \dhat\cdot\rvec_1 \\ \dhat\cdot\rvec_2 \end{pmatrix}
  \;}
  \label{eq:bgeo}
\end{equation}
Since $\rvec_1$ and $\rvec_2$ are the in-plane axes of the marker, seen from the
camera, $\bvec$ is the \emph{line of sight projected onto the marker plane}.

\paragraph{Step 3 --- Orthonormality.} The columns $\rvec_1,\rvec_2,\rvec_3$ of
a rotation form an orthonormal basis, so the squared components of any unit
vector along them add up to one:
\begin{equation}
  (\dhat\cdot\rvec_1)^2 + (\dhat\cdot\rvec_2)^2 + (\dhat\cdot\rvec_3)^2
  \;=\; \|\dhat\|^2 \;=\; 1 .
\end{equation}
The first two terms are $\|\bvec\|^2$ by \eqref{eq:bgeo}, hence
\begin{equation}
  \|\bvec\|^2 \;=\; 1 - (\dhat\cdot\rvec_3)^2 .
\end{equation}

\paragraph{Step 4 --- The viewing angle.} Substituting \eqref{eq:dn},
\begin{equation}
  \|\bvec\|^2 = 1 - \cos^2\theta = \sin^2\theta
  \qquad\Longrightarrow\qquad
  \boxed{\;\|\bvec\| = \sin\theta\;}
  \label{eq:bsin}
\end{equation}
The identity is exact and holds for \emph{every} pose, not only for particular
camera paths.

\subsection{The other quantities, for free}

The same reasoning gives the rest of $\Lmat$. Its entry $(3,3)$ is
$a = \Rvee\evec_3\cdot\rvec_3 = \dhat\cdot\rvec_3$, so by \eqref{eq:dn}
\begin{equation}
  a \;=\; -\cos\theta ,
  \label{eq:acos}
\end{equation}
and since $a = \tilde r_{11}\tilde r_{22} - \tilde r_{12}\tilde r_{21}
= \det\Rtil$ (Step~5), the block $\Rtil$ of a camera-facing marker always has a
negative determinant. From \eqref{eq:sv} and \eqref{eq:bsin} its singular
values are
\begin{equation}
  \sigma_1(\Rtil) = 1,
  \qquad
  \sigma_2(\Rtil) = \sqrt{1-\sin^2\theta} = \cos\theta ,
  \label{eq:svcos}
\end{equation}
and \eqref{eq:gamma} already gave $\gamma = 1/t_3$. Altogether,
\begin{equation}
  \gamma = \frac{1}{t_3},
  \qquad
  \sigma_2(\Rtil) = \cos\theta,
  \qquad
  \|\bvec\| = \sin\theta,
  \qquad
  a = -\cos\theta .
  \label{eq:geometry}
\end{equation}

\subsection{Numerical confirmation}

On a $0.5$~m arc with a $0.1$~m square, the solver's intermediate values agree
with \eqref{eq:geometry} to machine precision:

\begin{center}
\begin{tabular}{@{}rrrrrr@{}}
\toprule
$\theta$ & $\gamma$ & $1/t_3$ & $\sigma_2(\Rtil)$ & $\cos\theta$
  & $\|\bvec\|$ \\
\midrule
$0^{\circ}$  & 2.00000 & 2.00000 & 1.000000 & 1.000000 & 0.00000 \\
$5^{\circ}$  & 2.00000 & 2.00000 & 0.996195 & 0.996195 & 0.08716 \\
$20^{\circ}$ & 2.00000 & 2.00000 & 0.939693 & 0.939693 & 0.34202 \\
$45^{\circ}$ & 2.00000 & 2.00000 & 0.707107 & 0.707107 & 0.70711 \\
$70^{\circ}$ & 2.00000 & 2.00000 & 0.342020 & 0.342020 & 0.93969 \\
\bottomrule
\end{tabular}
\end{center}

\noindent
The identity is not specific to the arc: for the baseline pose, whose viewing
angle is $17.26^{\circ}$ and whose marker centre is off the optical axis, the
solver gives $\|\bvec\| = 0.296719 = \sin 17.26^{\circ}$, and its two components
match $(\dhat\cdot\rvec_1,\ \dhat\cdot\rvec_2)$ entry by entry.

\subsection{Interpretation}

Split the line of sight at the marker into a component along the normal, of
length $\cos\theta$, and a component lying in the marker plane, of length
$\sin\theta$. The vector $\bvec$ is exactly that in-plane component. Two facts
proved earlier now read geometrically:
\begin{itemize}[leftmargin=*]
  \item by the cancellation \eqref{eq:cancel}, the image derivative never sees
        $\bvec$: the local image of the plane does not reveal \emph{towards
        which side} of the normal the line of sight is tilted;
  \item by \eqref{eq:bb}, $\bvec$ is fixed only up to sign: the two candidates
        tilt the plane by the same angle $\theta$ on opposite sides of the line
        of sight.
\end{itemize}
Meanwhile $\Rtil$, whose singular values are $1$ and $\cos\theta$, is the
foreshortening itself: a plane seen obliquely is compressed by $\cos\theta$
along one direction and left unchanged along the perpendicular one. The image
carries the \emph{amount} of tilt, through $\cos\theta$, but not its
\emph{sign}; that missing sign is the two-pose ambiguity.

\paragraph{The degenerate case.} At $\theta = 0$ the map from the plane to the
image is a pure similarity, so $\Rtil$ is orthogonal and
$\Rtil^{\top}\Rtil = \Idm$. Equation \eqref{eq:bb} then gives
$\bvec\bvec^{\top} = \mathbf{0}$, hence $\bvec = \zerovec$; the cross product
\eqref{eq:third} becomes $(0,0,\pm1)^{\top}$, hence $\cvec = \zerovec$; and the
two matrices in \eqref{eq:twoposes} are \emph{identical}. The solver returns the
same pose twice, and the ambiguity is total: no image evidence can separate
poses that are equal.

% ============================================================
\section{Summary of the algorithm}
% ============================================================

\begin{enumerate}[leftmargin=*]
  \item Normalize the pixels with $\Kmat^{-1}$ and estimate $\Hmat$; scale it so
        that $h_{33} = 1$.
  \item Read $\vvec$ from \eqref{eq:vdef} and compute $\Jmat$ from
        \eqref{eq:JfromH}.
  \item Build $\Pimat$ from \eqref{eq:pidef}, $\Rvee$ from \eqref{eq:Rvdef}, and
        $\Bmat = \Pimat\Rvee^{(1,2)}$ from \eqref{eq:Bdef}.
  \item Compute $\Amat = \Bmat^{-1}\Jmat$, set $\gamma = \sigma_1(\Amat)$ and
        $\Rtil = \Amat/\gamma$; the depth is $t_3 = 1/\gamma$.
  \item Solve $\bvec\bvec^{\top} = \Idm - \Rtil^{\top}\Rtil$ for $\bvec$, which
        determines it up to a global sign.
  \item Complete each of the two rotations with
        $\lvec_3 = \lvec_1 \times \lvec_2$ and map back through
        $\Rmat^{\pm} = \Rvee\Lmat^{\pm}$.
  \item For each rotation, solve \eqref{eq:translation} for $\tvec$ by least
        squares.
  \item Reproject the four corners and rank the two candidates by pixel error.
\end{enumerate}

\end{document}
